% ============================== III. METHOD ==============================
\section{Method}
\label{sec:method}

\begin{figure}[t]
  \centering
  \resizebox{\textwidth}{!}{% Sơ đồ kiến trúc tổng thể (TikZ) — dùng trong figure* ở sections/03_method.tex
\begin{tikzpicture}[
  font=\scriptsize,
  >={Stealth[length=2mm]},
  box/.style={draw, rounded corners=2pt, align=center, minimum width=2.9cm, minimum height=0.95cm,
              inner sep=2pt, fill=white},
  nin/.style={box, fill=gray!10},
  nout/.style={box, fill=green!8},
  llm/.style={box, fill=orange!10},
  grp/.style={draw=#1!60, fill=#1!5, rounded corners=4pt, inner sep=6pt},
  lab/.style={font=\scriptsize\bfseries, text=#1!70!black},
  flow/.style={->, thick},
  aux/.style={->, densely dashed, gray!70!black},
]
% ---- cột 0: đầu vào
\node[nin] (lab)  at (0, 1.35)  {\textbf{Labels.txt} (HA4M)\\nhãn theo frame};
\node[nin] (vid)  at (0, 0)     {\textbf{Video RGB}\\.mp4 $|$ chuỗi ảnh $|$ HA4M};
\node[nin] (st)   at (0, -1.3)  {\textbf{station.yaml}\\zone, 12 bước SOP, VA};
\node[nin] (doc)  at (0, -3.9)  {\textbf{Tài liệu}\\WI / SOP / KB (PDF, DOCX, MD)};

% ---- cột 1: tầng skeleton
\node[box] (s1) at (3.75, 0)    {\textbf{01 Pose}\\MediaPipe $|$ RTMW $|$ RTMPose-ft};
\node[box] (s2) at (3.75, -1.3) {\textbf{02 Features}\\lọc, One Euro, tốc độ, zone};
\node[box] (s3) at (3.75, -2.6) {\textbf{03 Events}\\pick / reach / assemble / idle};

% ---- cột 2: tầng analyzer
\node[box] (a1) at (7.5, 0)     {\textbf{04 Steps}\\oracle $|$ rule $|$ MS-TCN};
\node[box] (a2) at (7.5, -1.3)  {\textbf{05 Metrics}\\Muda $\cdot$ Mura $\cdot$ Muri};
\node[box] (a3) at (7.5, -2.6)  {\textbf{06 Findings}\\R1--R8, xếp hạng s/chu kỳ};

% ---- cột 3: tầng LLM
\node[llm] (l2) at (11.25, -1.3) {\textbf{Validator}\\whitelist số + citation};
\node[llm] (l1) at (11.25, -2.6) {\textbf{07 Explain}\\Gemini $|$ Claude $|$ GPT};
\node[llm] (l0) at (11.25, -3.9) {\textbf{RAG store}\\chunk $\to$ e5 $\to$ Chroma};

% ---- cột 4: đầu ra
\node[nout] (o0) at (15.0, 0)    {\textbf{Run}\\manifest + artifact};
\node[nout] (o1) at (15.0, -1.3) {\textbf{08 Report}\\Markdown + PNG};
\node[nout] (o2) at (15.0, -2.6) {\textbf{UI / API}\\Streamlit, FastAPI};

% ---- nhóm tầng
\begin{scope}[on background layer]
  \node[grp=blue,   fit=(s1)(s3), label={[lab=blue]above:Tầng 1 — skeleton}] {};
  \node[grp=violet, fit=(a1)(a3), label={[lab=violet]above:Tầng 2 — analyzer (tất định)}] {};
  \node[grp=orange, fit=(l2)(l0), label={[lab=orange]above:Tầng 3 — LLM}] {};
\end{scope}

% ---- luồng chính
\draw[flow] (vid) -- node[above]{\tiny frame} (s1);
\draw[flow] (s1) -- node[right]{\tiny .npz} (s2);
\draw[flow] (s2) -- node[right]{\tiny .parquet} (s3);
\draw[flow] (s3.east) -- ++(0.35,0) |- (a1.west);
\draw[flow] (a1) -- node[right]{\tiny steps.csv} (a2);
\draw[flow] (a2) -- node[right]{\tiny metrics.json} (a3);
\draw[flow] (a3) -- node[above]{\tiny findings} (l1);
\draw[flow] (l1.north) -- (l2.south);
\draw[->, thick, red!70!black] (l2.east) .. controls +(0.55,-0.3) and +(0.55,0.3) .. node[right, align=left]{\tiny fail $\to$\\[-2pt]\tiny hỏi lại 1 lần} (l1.east);
\draw[flow] (l2) -- node[above]{\tiny passed} (o1);
\draw[flow] (o1) -- (o2);
\draw[aux]  (o1) -- (o0);
\draw[flow] (l0) -- node[right]{\tiny top-$k$ chunk} (l1);

% ---- luồng phụ
\draw[aux] (lab.east) -- (a1.north |- lab.east) node[pos=0.5, above]{\tiny chế độ oracle / eval} -- (a1.north);
\draw[aux] (st) -- node[above]{\tiny zone} (s2);
\draw[aux] (doc) -- node[above]{\tiny ingest (heading-aware 300--600 token)} (l0);
\draw[aux] ($(l0.west)+(0,0.22)$) -| node[pos=0.8, right]{\tiny TG chuẩn + ngưỡng} ($(a3.south)+(0.8,0)$);
\end{tikzpicture}
}
  \caption{Kiến trúc tổng thể AI Kaizen Copilot. Mỗi khối 01--08 là một chương trình đọc file vào và ghi file ra
  trong cùng một \textit{run} có manifest.
  Tầng 2 hoàn toàn tất định (không gọi LLM hay mạng). LLM ở tầng 3 chỉ diễn giải, và mọi câu trả lời phải qua validator.
  \code{station.yaml} còn cung cấp danh sách bước SOP và cờ VA cho khối 04--06.
  Đường nét đứt là dữ liệu phụ: nhãn HA4M cho chế độ oracle và đánh giá, thời gian chuẩn lấy từ tài liệu qua RAG.}
  \label{fig:arch}
\end{figure}

\subsection{Phát biểu bài toán}
\label{subsec:problem}
Đầu vào gồm ba phần:
\begin{itemize}
  \item video RGB $V=\{I_t\}_{t=1}^{T}$ của một trạm có một công nhân, tốc độ $f$ frame/s;
  \item cấu hình trạm $\mathcal{S}$, gồm tập vùng (zone) đa giác $\{Z_k\}$ trên ảnh (khay linh kiện, vùng lắp)
        và chuỗi bước SOP $\langle (j, \text{tên}_j, Z_{(j)}, \text{va}_j, \text{rep}_j)\rangle_{j=1}^{12}$;
  \item một kho tài liệu $\mathcal{D}$ (work instruction, SOP, cơ sở tri thức Kaizen).
\end{itemize}
Hệ thống cần xuất ra:
(a) phân đoạn chu kỳ thành các bước $\{(j, r, s_{j,r}, e_{j,r})\}$, với $r$ là lần lặp (ví dụ bước 2 lặp 3 lần);
(b) tập chỉ số $\mathcal{M}$, mỗi chỉ số có giá trị, đơn vị và định nghĩa;
(c) tập phát hiện $\mathcal{F}$, mỗi phát hiện kèm bằng chứng, ảnh hưởng (giây/chu kỳ) và độ tin cậy;
(d) với mỗi $f\in\mathcal{F}$, một giải thích $\langle$quan sát, giả thuyết, đề xuất, ECRS, trích dẫn$\rangle$
đã qua kiểm chứng.
Ràng buộc then chốt: mọi số xuất hiện trong (d) phải truy ngược được về (b)--(c).
\Cref{fig:arch} tóm tắt toàn bộ luồng xử lý.

\subsection{Ước lượng khung xương và làm sạch quỹ đạo cổ tay}
\label{subsec:pose}
Mỗi frame được đưa qua một backend pose trong vùng quan tâm (ROI) cố định của trạm.
Các backend là MediaPipe (33 điểm thân cộng $2\times21$ điểm tay, tổng 75), RTMW (133 điểm) hoặc RTMPose fine-tune (COCO-17).
Đầu ra là tensor $K\in\mathbb{R}^{T\times J\times 3}$ gồm $(x,y,\text{score})$.
Với mỗi cổ tay $h\in\{\text{L},\text{R}\}$, quỹ đạo $p^h_t$ được làm sạch theo bốn bước:
\begin{enumerate}
  \item điểm có score $<0{,}3$ coi là thiếu (NaN);
  \item \textit{điểm gai}, tức điểm mà vận tốc nhảy vào và nhảy ra đều vượt $v_{\max}=2500$\,px/s, bị loại;
  \item đoạn thiếu dài $\le 15$ frame (0{,}5\,s) được nội suy tuyến tính và đánh cờ \code{interp};
        đoạn dài hơn giữ NaN và không tính vào quãng đường;
  \item lọc One Euro \cite{casiez2012oneeuro} trên từng đoạn liên tục.
\end{enumerate}
Bộ lọc One Euro thích nghi tần số cắt theo tốc độ:
\begin{gather}
  \hat{p}_t = \alpha_t p_t + (1-\alpha_t)\hat{p}_{t-1},\qquad
  \alpha_t = \Big(1+\frac{f}{2\pi f_c(t)}\Big)^{-1},\nonumber\\
  f_c(t) = f_{c,\min} + \beta\,\lvert\dot{\hat p}_t\rvert,
  \label{eq:oneeuro}
\end{gather}
với $f_{c,\min}=1$\,Hz và $\beta=0{,}01$.
Tốc độ được tính \textit{sau} khi lọc: $v^h_t = f\,\lVert \hat{p}^h_t-\hat{p}^h_{t-1}\rVert_2$ (px/s).
Vùng của cổ tay là $z^h_t = \arg\{Z_k : \hat{p}^h_t\in Z_k\}$, theo thứ tự ưu tiên khi các đa giác chồng nhau;
nếu cổ tay không nằm trong vùng nào thì $z^h_t=$ \code{none}.

\subsection{Phát hiện sự kiện}
\label{subsec:events}
Trạng thái mỗi tay tại frame $t$ là cặp $\sigma^h_t=(z^h_t, m^h_t)$, trong đó $m\in\{\text{moving},\text{still},\text{lost}\}$
được xác định bằng ngưỡng $v^h_t \gtrless 120$\,px/s; \text{lost} ứng với cổ tay NaN.
Các frame liên tiếp cùng trạng thái được nén (run-length) thành sự kiện.
Sự kiện ngắn hơn $0{,}25$\,s được nhập vào sự kiện kề dài hơn, lặp lại cho tới khi không còn sự kiện ngắn,
rồi gộp các sự kiện kề nhau cùng trạng thái.
Mỗi sự kiện được gán loại theo quy tắc ở \cref{tab:events}.
Các quy tắc này cố ý đơn giản để kỹ sư đọc hiểu và chỉnh được.

\begin{table}[t]
\centering
\caption{Quy tắc suy diễn loại sự kiện (mỗi tay)}
\label{tab:events}
\footnotesize
\begin{tabularx}{\columnwidth}{@{}llX@{}}
\toprule
Vùng & Chuyển động & Loại sự kiện \\
\midrule
khay linh kiện & still $\le 2$\,s & \textit{pick} (lấy) \\
khay linh kiện & still $> 2$\,s  & \textit{search} (tìm) \\
khay linh kiện & moving           & \textit{reach} (với tay) \\
vùng lắp       & bất kỳ           & \textit{assemble} (lắp) \\
vùng thành phẩm & bất kỳ          & \textit{place} (đặt) \\
\code{none}    & still $> 1{,}5$\,s & \textit{idle} (chờ) \\
\code{none}    & still $\le 1{,}5$\,s & \textit{pause} \\
\code{none}    & moving           & \textit{transfer} (di chuyển) \\
--             & lost             & \textit{lost} (mất dấu) \\
\bottomrule
\end{tabularx}
\end{table}

\subsection{Phân đoạn bước SOP}
\label{subsec:steps}
Hệ thống hỗ trợ ba chế độ, đều ghi ra cùng một schema \code{step\_id, repetition, start\_s, end\_s, dur\_s, source}.

\textbf{Oracle.} Đọc nhãn HA4M theo frame. Mỗi đoạn liên tiếp cùng (action, repetition) thành một bước,
với $e = (\text{end\_frame}+1)/f$.

\textbf{Rule (baseline).} Các bước có khay riêng được neo bằng sự kiện lấy linh kiện (\cref{alg:rule}).
Bước không neo được, hoặc có mốc sai thứ tự, được gán \code{unclassified} và ghi vào danh sách bước thiếu.
Chế độ này không cần huấn luyện và giải thích được, nhưng không tách được những bước không có khay đặc trưng
(như bước 9, đặt khối 2 lên khối 1).

\begin{algorithm}[t]
\caption{Phân đoạn bước bằng mốc vùng (rule)}
\label{alg:rule}
\footnotesize
\begin{algorithmic}[1]
\Require Sự kiện $E$ (2 tay), chuỗi SOP $\langle (j,r,Z_{(j)})\rangle$, khoảng cách tối thiểu $\delta$
\Ensure Mốc bắt đầu $a_{j,r}$ cho các bước neo được
\State $t_{\text{prev}}\gets -\infty$
\For{mỗi $(j,r)$ theo thứ tự SOP}
  \If{$Z_{(j)}=\varnothing$} \textbf{đánh dấu} \code{no\_zone}; \textbf{continue} \EndIf
  \State $e^\star \gets$ sự kiện \textit{pick} hoặc \textit{reach} đầu tiên tại $Z_{(j)}$ có $s(e)>t_{\text{prev}}+\delta$
  \If{không có $e^\star$} \textbf{đánh dấu} \code{not\_found}; \textbf{continue} \EndIf
  \State $a_{j,r}\gets$ đầu chuỗi \textit{transfer/reach} liền trước $e^\star$ (lúc tay rời vùng lắp)
  \State $t_{\text{prev}}\gets a_{j,r}$
\EndFor
\State $e_{j,r}\gets$ mốc kế tiếp; bước cuối kết thúc tại sự kiện hoạt động cuối cùng
\State Khoảng giữa các mốc không gán được bước $\to$ \code{unclassified}
\end{algorithmic}
\end{algorithm}

\textbf{MS-TCN (học).} Một mạng MS-TCN \cite{farha2019mstcn} gán nhãn cho từng frame dựa trên đặc trưng khung xương
(\cref{subsec:arch}).
Mỗi lớp là một cặp (action, repetition), cộng thêm lớp nền ``0''.
Sau khi dự đoán, các mảnh bước ngắn hơn $0{,}5$\,s được nhập vào bước kề, trừ mảnh nền.
Sau đó các đoạn được ánh xạ về đúng schema của chế độ rule, nên các khối 05--08 không phải thay đổi.

\textbf{Đánh giá.} Khi có cả dự đoán và oracle, hệ thống ghép từng cặp bước theo $(j,r)$ và tính
$\mathrm{IoU}_{j,r}=\frac{|[s,e]\cap[\hat s,\hat e]|}{|[s,e]\cup[\hat s,\hat e]|}$.
Từ đó suy ra \Fone{} (bước đúng nếu IoU $\ge0{,}5$), MAE thời lượng, MAE mốc bắt đầu và kết thúc,
và sai số cycle time.

\subsection{Chỉ số Kaizen}
\label{subsec:metrics}
Gọi $d_{j}=\sum_r (e_{j,r}-s_{j,r})$ là thời lượng bước $j$.
Mọi chỉ số đều được lưu kèm đơn vị và định nghĩa ngắn; khối 07 đưa nguyên văn các định nghĩa này cho LLM.

\textbf{Muda.}
\begin{align}
  \mathrm{CT} &= e_{\text{cuối}} - s_{\text{đầu}}, \qquad
  \%\mathrm{VA} = 100\cdot\frac{\sum_{j:\,\text{va}_j} d_j}{\mathrm{CT}}, \label{eq:ct}\\
  L^h &= \sum_{t\in\mathcal{T}_{\text{hợp lệ}}} \lVert \hat p^h_t-\hat p^h_{t-1}\rVert_2, \label{eq:path}
\end{align}
trong đó $L^h$ là quãng đường cổ tay (px; đổi sang cm khi có hệ số \code{px\_per\_cm}).
Các chỉ số Muda khác: tổng thời gian idle và transfer của từng tay, số lần lấy theo từng zone trong một chu kỳ.

\textbf{Mura.} Hệ số biến thiên $\mathrm{CV}=\sigma/\mu$ của cycle time và của từng bước, tính giữa các lần lắp
của cùng một người và giữa các người lắp; liệt kê các bước có CV cao nhất.

\textbf{Muri (sàng lọc 2D).} Với khuỷu tay $\mathbf{e}$, cổ tay $\mathbf{w}$ và khớp gốc ngón giữa $\mathbf{m}$
(cùng một tay), góc cổ tay được tính bằng
\begin{equation}
  \theta_t = \arccos\frac{(\mathbf{w}-\mathbf{e})\cdot(\mathbf{m}-\mathbf{w})}
  {\lVert\mathbf{w}-\mathbf{e}\rVert\,\lVert\mathbf{m}-\mathbf{w}\rVert},\qquad
  \Delta_t = \lvert\theta_t - \operatorname{median}(\theta)\rvert .
  \label{eq:muri}
\end{equation}
Hệ thống báo phần trăm thời gian có $\Delta_t>30^\circ$, phân vị 90 của $\Delta$, và độ phủ (tỷ lệ frame đo được).
Đây là góc chiếu 2D, so với trung vị của chính người đó; nó chỉ dùng để sàng lọc, không thay thế
đánh giá công thái học chuyên sâu như RULA.

\subsection{Phát hiện bằng quy tắc}
\label{subsec:findings}
Tám quy tắc ở \cref{tab:rules} chạy trên các chỉ số.
Thêm một quy tắc chỉ cần thêm một phần tử vào danh sách \code{RULES}.
Mỗi phát hiện ghi bằng chứng (chỉ số, giá trị, đơn vị, trial, các mốc thời gian), ảnh hưởng $I$ (giây/chu kỳ)
và độ tin cậy $c\in\{\text{high},\text{med},\text{low}\}$.
Mức nghiêm trọng là \textit{high} nếu $I/\mathrm{CT}\ge5\%$, \textit{med} nếu $I/\mathrm{CT}\ge2\%$, còn lại là \textit{low}.
Danh sách phát hiện được xếp hạng theo
\begin{equation}
  \text{score}(f) = I(f)\cdot w(c),
  \label{eq:rank}
\end{equation}
với $w(\text{high})=1$, $w(\text{med})=0{,}6$, $w(\text{low})=0{,}3$.
Thời gian chuẩn $s^{\text{std}}_j$ và ngưỡng của R8 ưu tiên giá trị ghi trong \code{station.yaml}.
Nếu không có, hệ thống trích từ bảng ``TG chuẩn'' của work instruction trong kho RAG,
và mỗi giá trị đều mang trích dẫn tới trang nguồn.

\begin{table}[t]
\centering
\caption{Các quy tắc phát hiện (V: chạy với 1 video; A: cần nhiều video)}
\label{tab:rules}
\footnotesize
\setlength{\tabcolsep}{3pt}
\begin{tabularx}{\columnwidth}{@{}clXc@{}}
\toprule
ID & Loại & Điều kiện kích hoạt & Phạm vi \\
\midrule
R1 & bước chậm & $d_j > \mathrm{median}+2\cdot\mathrm{MAD}$ (giữa các trial, cần $\ge3$) & A \\
R2 & lấy dư & số pick tại zone $>$ $\sum$ rep của các bước dùng zone đó & V \\
R3 & idle dài & tổng idle 1 tay $>3$\,s/chu kỳ & V \\
R4 & sai thứ tự & chuỗi bước $\neq$ SOP hoặc thiếu bước & V \\
R5 & lệch tay & 1 tay idle $>60\%$ bước khi tay kia hoạt động $>60\%$ & V \\
R6 & biến động & CV cycle time trong 1 người $>0{,}25$ & A \\
R7 & Muri cổ tay & \% thời gian $\Delta_t>30^\circ$ vượt ngưỡng & V \\
R8 & lệch chuẩn & $\lvert d_j-s^{\text{std}}_j\rvert/s^{\text{std}}_j>30\%$ (từ tài liệu) & V \\
\bottomrule
\end{tabularx}
\end{table}

\subsection{Diễn giải bằng LLM có kiểm chứng}
\label{subsec:llm}
\textbf{Prompt.} Prompt có version (\code{explain\_v2}) gồm: định nghĩa chỉ số, phát hiện ở dạng JSON, bảng bước kèm thời gian
chuẩn, và ngữ cảnh tri thức.
Ngữ cảnh tri thức có hai dạng: (a) toàn bộ KB/SOP nhỏ với các đoạn được đánh nhãn \code{[KB-xx]}, \code{[SOP-xx]};
hoặc (b) top-$k$ chunk truy xuất từ kho RAG, mỗi chunk có \code{chunk\_id} dạng \code{DOC-xxxxxx:pN:cM}.
LLM phải trả về JSON thuần gồm quan sát ($\le2$ câu), giả thuyết (1 câu), đề xuất ($\le2$ câu), nhãn ECRS,
hướng dẫn đo lại, và danh sách trích dẫn.

\textbf{Kho tài liệu.} Tài liệu (PDF, DOCX, MD, TXT) được cắt theo heading, mỗi chunk 300--600 token, chồng lấn 50 token.
Mỗi chunk thuộc đúng một trang để trích dẫn chính xác trang; bảng được giữ thành chunk riêng.
Chunk được nhúng bằng multilingual-E5 \cite{e5} và lưu cục bộ trong Chroma \cite{chroma}.

\textbf{Validator.} Với phát hiện $f$, gọi $\mathcal{X}(f)$ là tập các giá trị số trong bằng chứng của nó
và $\mathcal{N}(f)$ là tập số nằm trong tên trường hoặc chuỗi
(ví dụ \code{dev\_p90\_deg}$\to$90, ``lệch $>30^\circ$''$\to$30).
Whitelist của $f$ là
\begin{multline}
  \mathcal{W}(f)=\mathcal{U}\ \cup\ \mathcal{N}(f)\\ \cup \bigcup_{x\in\mathcal{X}(f)}
  \{\mathrm{rd}_0(x), \mathrm{rd}_1(x), \mathrm{rd}_0(|x|), \mathrm{rd}_1(|x|)\},
  \label{eq:whitelist}
\end{multline}
với $\mathcal{U}=\{0,1,2,3,100\}\cup\{1,\dots,12\}$ (số phổ quát và số thứ tự bước), $\mathrm{rd}_k$ là làm tròn
$k$ chữ số thập phân.
Một câu trả lời \textit{đạt} khi cả ba điều kiện sau thoả:
\begin{enumerate}
  \item mọi token số $u$ trong ba trường văn bản thoả $\mathrm{rd}_1(u)\in\mathcal{W}(f)$;
        token dính chữ phía trước như \code{P90}, \code{U001} được bỏ qua;
  \item mọi trích dẫn, sau khi chuẩn hoá (\code{kb13}, \code{[KB-13]} thành \code{KB-13}), tồn tại trong KB/SOP
        hoặc trong kho RAG;
  \item định dạng đúng: số câu nằm trong giới hạn và ECRS $\in\{E,C,R,S\}$.
\end{enumerate}
Nếu không đạt, hệ thống gọi lại LLM \textit{một lần}, kèm thông báo chỉ liệt kê phần sai
(ví dụ ``số lạ: [`8.2']'', ``citation không tồn tại: [`KB-99']'').
Kết quả lưu mọi lần thử cùng trạng thái \code{passed}, \code{failed\_after\_retry} hoặc \code{parse\_failed}.
Pipeline không dừng khi LLM lỗi.
Hai chỉ số được báo cáo là tỷ lệ đạt ngay lần đầu và tỷ lệ đạt sau khi hỏi lại.

\subsection{Báo cáo và truy vết}
Khối 08 sinh báo cáo Markdown gồm: tóm tắt và đề xuất ($\le6$ dòng), bảng bước kèm thời gian chuẩn, độ lệch và nguồn,
biểu đồ Gantt và Yamazumi, 3M tổng hợp, Muri, các phát hiện kèm ảnh bằng chứng, bảng tài liệu tham chiếu
và mục ``Giả định \& giới hạn'' tự sinh.
Ảnh bằng chứng là frame tại thời điểm của phát hiện, có vẽ cổ tay và các zone.
Bảng tài liệu tham chiếu được kiểm tra khớp hai chiều với mọi nhãn trích dẫn trong thân báo cáo.
Mục ``Giả định \& giới hạn'' ghi rõ fps, đơn vị px hay cm, oracle hay dự đoán, và việc Muri chỉ là góc 2D.
